\chapter{Background and Related Work}

This chapter sets up the technical context the thesis builds on: the
online bipartite-matching problem and its input models (§2.1), the
known-i.i.d. model we work in (§2.2), the learning-augmented
(``with-predictions'') paradigm and its consistency/robustness language
(§2.3), the specific prediction-based matching algorithms we benchmark
(§2.4), the distribution-testing results behind the algorithms' tests
and the outlook of §10.2 (§2.5), and the gaps in the literature this
thesis addresses (§2.6).

\section{Online bipartite matching and competitive
analysis}

In online bipartite matching, one side of a bipartite graph (the
\emph{offline} resources) is known in advance, while the vertices of the
other side (the \emph{online} requests) arrive one at a time. On each
arrival the algorithm sees the request's incident edges and must either
match it to a currently unmatched neighbor or leave it unmatched,
\textbf{immediately and irrevocably}. The goal is to maximize the size
(or, in weighted variants, the value) of the matching. Performance is
measured by the \textbf{competitive ratio}: the ratio between the
algorithm's matching and an offline optimum on the same input, in
expectation over both the random input and the algorithm's own
randomness (§3.1 fixes the precise form we report).

The difficulty of the problem depends entirely on the assumed
\emph{input model}:

\begin{itemize}
\tightlist
\item
  \textbf{Adversarial arrival order.} The requests and their order are
  chosen by an adversary. The seminal result of Karp, Vazirani and
  Vazirani \autocite{kvv1990ranking} is that the randomized Ranking
  algorithm (fix a uniformly random priority order on the resources and
  match each request to its highest-priority available neighbor)
  achieves competitive ratio \(1-1/e\approx0.632\), and that this is
  optimal for the adversarial model. Deterministic algorithms cannot
  beat \(1/2\).
\item
  \textbf{Random arrival order.} The graph is adversarial but the
  requests arrive in a uniformly random order; this is strictly easier
  than adversarial order and admits ratios above \(1-1/e\).
\item
  \textbf{Known-i.i.d.} Requests are drawn independently from a known
  distribution over request types (§2.2); this is the average-case model
  and the one this thesis studies.
\end{itemize}

These three models are nested in difficulty, and the direction matters
throughout the thesis: every known-i.i.d. instance is also a
random-order instance, and every random-order instance is an
adversarial-order instance. A guarantee proved in a harder model
therefore holds in an easier one, but not conversely, so a bound proved
for adversarial order says nothing about how well an algorithm does on
known-i.i.d. inputs beyond that floor.

\section{The known-i.i.d. model}

In the known-i.i.d. model there is a bipartite \emph{type graph}: each
online \textbf{type} \(\ell\) has a fixed neighborhood \(N(\ell)\) among
the offline resources, and an instance is generated by drawing \(m\)
arrivals independently from a known distribution \(p\) over types. The
type graph and \(p\) are known; the realized sequence is not. This
models settings (ad serving, recurring request streams) where the
population of requests is statistically stable even though individual
arrivals are not.

A line of work has pushed the worst-case (over type graphs) competitive
ratio above the \(1-1/e\) adversarial barrier: Feldman et al.
\autocite{feldman2009online} first beat it (\(0.670\)) via a flow-based
\emph{blue/red} decomposition of a suggested matching; Manshadi, Oveis
Gharan and Saberi \autocite{manshadi2012online} and Jaillet--Lu
\autocite{jailletlu2014online} improved the ratio (to \(\approx0.702\)
and \(\approx0.729\) respectively) using LP-based and list-based online
policies; subsequent work refined it further. These are the
``sophisticated'' algorithms whose worst-case optimality motivates their
use.

Crucially for this thesis, Borodin, Karavasilis and Pankratov
\autocite{borodin2018experimental} conducted an experimental study of
these algorithms and found that on average-case i.i.d. instances the
picture is very different from the worst case: the simple algorithms
(Greedy, Ranking) perform almost as well as the sophisticated ones,
whose advantage is a worst-case, not a typical-case, phenomenon.

One empirical observation in that line is the seed of this thesis:
\emph{on typical inputs the simple baseline is already near-optimal}.
Everything that follows is an attempt to find out what a prediction can
add to it, and we reproduce a subset of the study as validation (Chapter
3).

\section{Learning-augmented
algorithms}

The \textbf{algorithms-with-predictions} (or learning-augmented)
paradigm, surveyed in \autocite{mitzenmacher2020algorithms}, equips an
online algorithm with a prediction about the unknown input and asks for
two guarantees simultaneously:

\begin{itemize}
\tightlist
\item
  \textbf{Consistency:} near-optimal performance when the prediction is
  accurate;
\item
  \textbf{Robustness:} performance no worse than a prediction-free
  guarantee when the prediction is arbitrarily wrong.
\end{itemize}

The paradigm was crystallized by Lykouris and Vassilvitskii
\autocite{lykouris2018caching} for competitive caching, who showed how
to interpolate between trusting and ignoring a predictor while bounding
both consistency and robustness. A large literature followed across ski
rental, scheduling and other online problems, including the optimal
consistency/robustness trade-off analyses of Wei and Zhang
\autocite{weizhang2020tradeoffs}, which establish problem-intrinsic
Pareto frontiers between the two objectives. Recent work further
analyzes how the achievable ratio degrades continuously with prediction
accuracy in a distributionally-robust formulation
\autocite{yoshinaga2026accuracy}, complementing the discrete
consistency/robustness endpoints used here.

Two items from this literature are load-bearing for us. The recurring
mechanism is \emph{hedging}: combining the prediction with a safe
default so that a bad prediction cannot cause catastrophe. The
blind-follow-with-switching combiner of Chłędowski, Polak, Szabucki and
Żołna \autocite{chledowski2021caching} is one such mechanism, which we
port and benchmark (Chapter 6); their paper, an experimental study of
robust learning-augmented caching, is also the closest methodological
template for this thesis's experimental half.

\section{Predictions for online bipartite
matching}

Two strands apply the with-predictions paradigm to online matching,
differing in the \emph{prediction object} they consume.

\textbf{Degree predictions: MinPredictedDegree.} Aamand, Chen and Indyk
\autocite{aci2022mpd} propose \textbf{MinPredictedDegree (MPD)}: given a
prediction \(\mu\) of each offline resource's degree (how contended it
will be), match each arrival to its available neighbor of \emph{minimum
predicted degree}, i.e.~protect the resources predicted to be rarest.
MPD is robust by construction: a constant (useless) predictor reduces it
to Ranking. A key structural fact, which the thesis engages in Chapter
5, is that MPD depends on \(\mu\) \emph{only through the order it
induces}. The authors' Appendix D bounds the matching loss by an order
quantity built in two steps: list the true degrees in the order the
prediction suggests, then count how many of them are out of place:
formally \(n-\mathrm{LIS}\), where \(\mathrm{LIS}\) is the length of the
longest non-decreasing subsequence of that list. The count is zero
exactly when the prediction gets the order right, so a monotone
(order-preserving) prediction incurs zero loss.

\textbf{Type-histogram advice and test-and-fallback.} A second strand
takes the prediction to be a histogram \(\hat c\) over request types
(how many of each type will arrive). Choo et al.
\autocite{choo2024imperfect} introduce TestAndMatch: build a matching
from \(\hat c\) and Mimic it, but first spend a sublinear prefix of
arrivals testing whether the observed type frequencies match \(\hat c\),
using an \(\ell_1\)-distance tester adapted from Jiao, Han and Weissman
\autocite{jiao2018l1}, and fall back to Ranking if they do not.
Burathep, Erlebach and Moses \autocite{bem2026testmatch} generalize this
(``Test-and-Match+'') to the random-arrival model and to imperfect
knowledge of the matching size. Both papers give \emph{upper} bounds
(algorithms); their only lower bound (Choo et al.'s Theorem 3.1) is a
generic adversarial indistinguishability result (no algorithm is
\(1\)-consistent and \(>\tfrac12\)-robust), and neither proves a lower
bound in the stochastic model. Notably, Choo et al.'s acceptance
threshold already couples to the baseline competitive ratio \(\beta\),
but constructively, inside the algorithm design, not as a lower bound.
What that coupling costs, i.e.~how large a prefix the follow/fallback
decision fundamentally requires, is the question this thesis measures
(Chapter 6) and then reads quantitatively (§10.2).

\section{Distribution testing}

The test at the heart of the algorithms above is a question in
\textbf{distribution testing}: given samples from an unknown
distribution \(p\) over a support of size \(r\) (the same \(r\) as in
our model, the number of request types, because the histogram advice is
a distribution over types) and a known reference \(q\), decide how far
\(p\) is from \(q\). We measure that distance in \(\ell_1\) throughout,
as the algorithms and our experiments do; it is \emph{twice} the
total-variation distance, and every threshold quoted in this thesis is
an \(\ell_1\) threshold. Two regimes must be distinguished:

\begin{itemize}
\tightlist
\item
  \textbf{Identity (non-tolerant) testing}, distinguishing \(p=q\) from
  \(\lVert p-q\rVert_1\ge
  \varepsilon\), has sample complexity \(\Theta(\sqrt r/\varepsilon^2)\)
  \autocite{paninski2008coincidence,valiant2017automatic},
  \emph{sublinear} in the support size.
\item
  \textbf{Tolerant testing / distance estimation}, distinguishing
  \(\lVert p-q\rVert_1\le
  \varepsilon_1\) from \(\ge\varepsilon_2\) for
  \(0<\varepsilon_1<\varepsilon_2\), i.e.~estimate the distance rather
  than test equality, is provably much harder: Valiant and Valiant
  \autocite{valiant2011unseen} showed it requires \(\Theta(r/\log r)\)
  samples, \emph{near-linear} in the support. Jiao, Han and Weissman
  \autocite{jiao2018l1} gave matching bounds for \(\ell_1\)-distance
  estimation, and Canonne, Jain, Kamath and Li
  \autocite{canonne2022tolerance} precisely characterized the whole
  \((\varepsilon_1,\varepsilon_2)\) landscape, showing that for constant
  tolerances the cost jumps to the ``barely sublinear''
  \(\tilde\Theta(r/\log r)\).
\end{itemize}

The near-quadratic gap between \(\sqrt r\) (testing) and \(r/\log r\)
(tolerant testing) matters twice in this thesis. It is why the
deployable versions of TestAndMatch fall back to an empirical surrogate
for their tester, and it is why that surrogate is blind on large
supports: the resolution limit measured in Chapter 6. Whether the
barrier also dooms every other decision rule turns out to be subtler
than it first appears; the concluding outlook (§10.2) returns to this
point.

\section{Positioning of this thesis}

Against this background, three gaps stand out. They correspond
one-to-one to the three questions of §1.2, and the thesis addresses them
in that order:

\begin{enumerate}
\def\labelenumi{\arabic{enumi}.}
\tightlist
\item
  \textbf{No unified empirical comparison found in the reviewed
  literature.} The matching algorithms above were studied on their own
  input families and error models, and largely in theory; we did not
  identify a head-to-head benchmark under a common harness. (Chapters
  4--7.)
\item
  \textbf{No empirical study of test-and-fallback found in the reviewed
  literature.} The distribution test at the heart of
  \autocite{choo2024imperfect,bem2026testmatch} has no deployable
  implementation (the authors themselves fall back to an empirical
  surrogate), and its testing cost, threshold calibration, and failure
  modes have not been measured. (Chapter 6.)
\item
  \textbf{No quantification of what the prefix test costs.} No prior
  work measures, or bounds, how large a prefix the follow/fallback
  decision requires on strong-baseline instances, where the upside to be
  captured is smallest (§8.3 describes the prior-art pass behind this
  claim). (Chapter 6 empirically; §10.2 in outlook.)
\end{enumerate}
